% Provenance (never printed): see sources/outline_sources.md (Year 2: tholeiitic, calc-alkaline
% and alkaline series; the TAS diagram; settings; isotopic signatures of mantle and crust).
% Classification from BGS RR 99-06 (after Streckeisen 1976, Le Bas et al. 1986, Le Maitre et
% al. 1989); series and data from Vermeesch and Pease (2021) and the GeoplotR package;
% curves and counts from figdata/bachelor-2/series.py; exercise and problem values in the
% Book 4 scratch calculators calc/ch09.py and calc/ch09pb.py.
\chapter{Séries magmatiques et contextes tectoniques}\label{ch:b2:magmatic-series}

Donnez à une pétrologue une analyse chimique sans carte, et elle dira souvent si la
lave est sortie au niveau d'une dorsale, sur une île océanique ou au-dessus d'une plaque
en subduction. Les laves d'Islande et celles des volcans des Cascades, dans l'ouest de
l'Amérique du Nord, vont toutes deux du basalte à la rhyolite~; analysées par centaines,
elles suivent deux trajectoires différentes au fil de leur évolution, l'une s'enrichissant
sans cesse en fer, l'autre non. Ce chapitre nomme d'abord les roches magmatiques, d'après
leurs minéraux lorsqu'on peut les voir et d'après leur chimie sinon~; il les regroupe ensuite
en séries, familles de roches apparentées par différenciation, et relie chaque série à un contexte
tectonique, avec les diagrammes que les pétrologues utilisent pour lire ce contexte dans une analyse, et leurs limites.
% ledger: vp21:series

\begin{recall}
\textbf{Du volume de lycée}~: basalte, andésite, rhyolite, gabbro et granite~;
roches felsiques et mafiques~; textures volcaniques et plutoniques. \textbf{Du volume de
première année}~: dorsales océaniques, zones de subduction et arcs volcaniques, points chauds
et grandes provinces ignées, rifts continentaux. \textbf{Du \cref{ch:b2:elements-in-the-earth}}~:
MORB et OIB, \omterm{def:b2:elements-in-the-earth:normalised}{diagrammes normalisés} au \omterm{def:b2:elements-in-the-earth:bse}{manteau primitif} (en araignée), \omterm{def:b2:elements-in-the-earth:partition}{éléments compatibles} et
incompatibles. \textbf{Du \cref{ch:b2:isotope-geochemistry}}~: le \omterm{def:b2:isotope-geochemistry:reservoirs}{manteau appauvri} et les
réservoirs isotopiques enrichis. \textbf{Du
\cref{ch:b2:magmatic-differentiation}}~: \omterm{def:b2:magmatic-differentiation:crystallisation}{cristallisation fractionnée}, ligne de descente du
liquide, \omterm{def:b2:magmatic-differentiation:harker}{diagrammes de Harker} et loi de Rayleigh.
\end{recall}

\section{Nommer les roches magmatiques}\label{sec:b2:magmatic-series:naming}

Une roche grenue montre ses minéraux, et on la nomme d'après eux. Une lave est souvent trop fine
ou trop vitreuse pour cela, et on la nomme d'après son analyse chimique. Les deux méthodes sont
des conventions internationales~: elles décrivent, sans rien dire de la manière dont la roche
s'est formée.
% ledger: bgs:tas-use

\begin{definition}[Composition modale~; diagramme QAPF]\label{def:b2:magmatic-series:qapf}
La \emph{composition modale}\index{composition modale} d'une roche est la proportion, en
volume, de chacun de ses minéraux. Le \emph{diagramme QAPF}\index{diagramme QAPF} nomme une
roche magmatique grenue dont moins de \qty{90}{\percent} des minéraux sont mafiques, à partir de
quatre composants modaux recalculés à une somme de 100~: quartz (Q), feldspath alcalin (A),
plagioclase (P) et feldspathoïdes (F). Puisque le quartz et les feldspathoïdes ne cristallisent pas
ensemble, chaque roche se place dans l'un de deux triangles, QAP ou FAP, reliés par leur côté
commun AP.
\end{definition}
% ledger: bgs:qapf

\begin{omfigure}
\centering
\begin{tikzpicture}[font=\scriptsize, x=1cm, y=1cm]
\def\L{6.4}
\pgfmathsetmacro{\H}{\L*0.8660254}
% upper point (q, p): q = Q fraction, p = P/(A+P)
\newcommand{\qapfup}[2]{({(#1)*\L/2 + (#2)*(1-(#1))*\L}, {(#1)*\H})}
\newcommand{\dn}[2]{({(#1)*\L/2 + (#2)*(1-(#1))*\L}, {-(#1)*\H})}
\draw[thick] (0,0) -- (\L/2,\H) -- (\L,0) -- (\L/2,-\H) -- cycle;
\draw (0,0) -- (\L,0);
\foreach \q in {0.9,0.6,0.2,0.05} \draw \qapfup{\q}{0} -- \qapfup{\q}{1};
\foreach \q in {0.1,0.6} \draw \dn{\q}{0} -- \dn{\q}{1};
\foreach \p in {0.1,0.65,0.9} \draw \qapfup{0.6}{\p} -- \qapfup{0}{\p};
\draw[dashed] \qapfup{0.6}{0.35} -- \qapfup{0.2}{0.35};
\draw \qapfup{0.2}{0.35} -- \qapfup{0}{0.35};
\foreach \p in {0.1,0.35,0.65,0.9} \draw \dn{0}{\p} -- \dn{0.1}{\p};
\foreach \p in {0.1,0.5,0.9} \draw \dn{0.1}{\p} -- \dn{0.6}{\p};
\node[anchor=south] at (\L/2,\H) {\normalsize Q};
\node[anchor=east] at (0,0) {\normalsize A};
\node[anchor=west] at (\L,0) {\normalsize P};
\node[anchor=north] at (\L/2,-\H) {\normalsize F};
\node at \qapfup{0.94}{0.5} {1a};
\node at \qapfup{0.75}{0.5} {1b};
\node[font=\tiny] at \qapfup{0.4}{0.05} {2};
\node at \qapfup{0.4}{0.225} {3a};
\node at \qapfup{0.4}{0.5} {3b};
\node at \qapfup{0.4}{0.775} {4};
\node[font=\tiny] at \qapfup{0.4}{0.95} {5};
\foreach \p/\n in {0.225/7*,0.5/8*,0.775/9*} \node at \qapfup{0.125}{\p} {\n};
\node[font=\tiny] at \qapfup{0.125}{0.05} {6*};
\node[font=\tiny] at \qapfup{0.125}{0.95} {10*};
\foreach \p/\n in {0.225/7,0.5/8,0.775/9} \node at \qapfup{0.025}{\p} {\n};
\foreach \p/\n in {0.225/7',0.5/8',0.775/9'} \node at \dn{0.05}{\p} {\n};
\node[font=\tiny] at \dn{0.3}{0.05} {11};
\node at \dn{0.35}{0.3} {12};
\node at \dn{0.35}{0.7} {13};
\node[font=\tiny] at \dn{0.3}{0.95} {14};
\node at \dn{0.75}{0.5} {15};
\node[anchor=east] at \qapfup{0.9}{0} {90};
\node[anchor=east] at \qapfup{0.6}{0} {60};
\node[anchor=east] at \qapfup{0.2}{0} {20};
\node[anchor=east] at \dn{0.1}{0} {10};
\node[anchor=east] at \dn{0.6}{0} {60};
\foreach \p/\n in {0.1/10,0.35/35,0.65/65,0.9/90} \node[anchor=south, fill=white, inner sep=0.5pt] at \qapfup{0.6}{\p} {\tiny\n};
\node[anchor=west, align=left, font=\scriptsize] at (\L+0.9,1.6) {Triangle supérieur~: \% Q\\dans Q + A + P~;\\
  droites où\\$100\,P/(A+P)$\\vaut 10, 35, 65, 90};
\node[anchor=west, align=left, font=\scriptsize] at (\L+0.9,-1.7) {Triangle inférieur~: \% F\\dans F + A + P~;\\
  droites où\\$100\,P/(A+P)$\\vaut 10, 50, 90};
\end{tikzpicture}
\par\smallskip{\small Classification QAPF des roches magmatiques grenues. 1a quartzolite,
1b granitoïde riche en quartz, 2 granite à feldspath alcalin, 3a syénogranite et 3b monzogranite
(ensemble les granites), 4 granodiorite, 5 tonalite~; 6*--10* les variétés à quartz des champs
6 à 10~: 6 syénite à feldspath alcalin, 7 syénite, 8 monzonite, 9 monzodiorite ou
monzogabbro, 10 diorite, gabbro ou anorthosite~; 6'--10' leurs variétés à feldspathoïdes~;
11 syénite à feldspathoïdes, 12 monzosyénite à feldspathoïdes, 13 monzodiorite ou monzogabbro
à feldspathoïdes, 14 diorite ou gabbro à feldspathoïdes, 15 foïdolite. Les champs minces 6, 6',
10 et 10' le long des bords ne sont pas numérotés. D'après le schéma du British Geological Survey
(fondé sur Streckeisen, 1976).}
\end{omfigure}
% ledger: bgs:qapf

Le champ 10 regroupe trois noms. Le \omterm{def:b2:magmatic-series:qapf}{diagramme QAPF} ne sépare pas la diorite, le gabbro et
l'anorthosite, ni les roches gabbroïques entre elles~; d'autres diagrammes, fondés sur le
plagioclase et les minéraux mafiques, le font.

\begin{definition}[Diagramme alcalins--silice]\label{def:b2:magmatic-series:tas}
Le \emph{diagramme alcalins--silice}\index{diagramme alcalins--silice} (diagramme TAS) nomme
une roche volcanique dont la \omterm{def:b2:magmatic-series:qapf}{composition modale} ne peut être mesurée, à partir de ses pourcentages
massiques de \ce{Na2O + K2O} par rapport à \ce{SiO2}, une fois l'analyse recalculée à
\qty{100}{\percent} sans eau ni dioxyde de carbone. Il est divisé en quinze champs par des
droites dont les intersections ont des coordonnées fixes
(la figure après le paragraphe suivant).
\end{definition}
% ledger: bgs:tas, bgs:tas-use

Le diagramme comporte trois groupes de champs. Le long du bas, à faibles teneurs en alcalins,
figurent les noms des séries courantes, du picrobasalte au basalte, puis jusqu'à la rhyolite,
séparés à 45, 52, 57 et \qty{63}{\percent} de silice et, pour la rhyolite, par une droite
inclinée. Au-dessus d'eux, les champs S et le trachyte regroupent des roches plus riches en
alcalins à même teneur en silice~; les champs U, la phonolite et la foïdite regroupent des roches
si riches en alcalins, et si pauvres en silice, qu'elles cristallisent des feldspathoïdes.
Quelques champs se subdivisent selon un second critère~: un trachybasalte s'appelle hawaiite si $\ce{Na2O} - 2
\ge \ce{K2O}$, trachybasalte potassique sinon~; une téphrite dont la norme contient plus de
\qty{10}{\percent} d'olivine (ci-dessous) est une basanite~; une rhyolite ou un trachyte dont le
rapport molaire $(\ce{Na2O + K2O})/\ce{Al2O3}$ dépasse 1 est dit hyperalcalin. Le schéma exige des
analyses contenant moins de \qty{2}{\percent} d'eau et \qty{0.5}{\percent} de dioxyde de carbone~:
l'altération ajoute de l'eau et déplace les alcalins, et un \omterm{def:b2:magmatic-differentiation:cumulate}{cumulat} n'est pas un liquide.
% ledger: bgs:tas, bgs:tas-use

\begin{omfigure}
\centering
\begin{tikzpicture}
\begin{axis}[omrecord, width=0.86\linewidth, height=8.6cm,
  xmin=37, xmax=80, ymin=0, ymax=16, xtick={37,41,...,77}, ytick={0,2,...,16},
  xlabel={\ce{SiO2} (\unit{\percent})}, ylabel={\ce{Na2O + K2O} (\unit{\percent})},
  legend style={font=\scriptsize, at={(0.98,0.98)}, anchor=north east},
  tick label style={font=\scriptsize}, label style={font=\small}]
\addplot[only marks, mark=*, mark size=0.55pt, blue!70!black, opacity=0.6] table[x=SiO2, y=alk] {figdata/out/bachelor-2-series-cath-th.dat};
\addlegendentry{Islande}
\addplot[only marks, mark=*, mark size=0.55pt, red!75!black, opacity=0.6] table[x=SiO2, y=alk] {figdata/out/bachelor-2-series-cath-ca.dat};
\addlegendentry{Cascades}
\foreach \f in {Pc,B,O1,O2,O3,R,S1,S2,S3,T,U1,U2,U3,Ph,F}
  \addplot[black, thin, forget plot] table[x=SiO2, y=alk] {figdata/out/bachelor-2-series-tas-\f.dat};
\node[font=\tiny] at (axis cs:43,1.6) {Pc};
\node[font=\tiny, fill=white, inner sep=0.5pt] at (axis cs:50.5,0.8) {B};
\node[font=\tiny] at (axis cs:54.5,2.2) {O1};
\node[font=\tiny] at (axis cs:60,2.2) {O2};
\node[font=\tiny] at (axis cs:66.5,2.2) {O3};
\node[font=\tiny] at (axis cs:74.5,11) {R};
\node[font=\tiny, fill=white, inner sep=0.5pt] at (axis cs:49,5.8) {S1};
\node[font=\tiny] at (axis cs:52.5,7.4) {S2};
\node[font=\tiny] at (axis cs:57.5,8.9) {S3};
\node[font=\tiny] at (axis cs:64.5,12) {T};
\node[font=\tiny] at (axis cs:44,6.6) {U1};
\node[font=\tiny] at (axis cs:48.7,9.4) {U2};
\node[font=\tiny] at (axis cs:52.5,11.8) {U3};
\node[font=\tiny] at (axis cs:57,14.6) {Ph};
\node[font=\tiny] at (axis cs:39,10) {F};
\end{axis}
\end{tikzpicture}
\par\smallskip{\small Le diagramme TAS (d'après Le Bas et al., 1986, tel que donné par le
British Geological Survey), avec 456 laves et roches intrusives d'Islande (bleu) et 580 des
Cascades (rouge), oxydes tels qu'analysés. Pc picrobasalte, B basalte, O1 andésite basaltique,
O2 andésite, O3 dacite, R rhyolite~; S1 trachybasalte, S2 trachyandésite basaltique,
S3 trachyandésite, T trachyte ou trachydacite~; U1 téphrite ou basanite, U2 phonotéphrite,
U3 téphriphonolite, Ph phonolite, F foïdite.}
\end{omfigure}
% ledger: bgs:tas, geoplotr:tas, data:b4-cath


\begin{definition}[Saturation en silice]\label{def:b2:magmatic-series:saturation}
Un \omterm{def:b2:magma-emplacement:magma}{magma} ou une roche est \emph{sursaturé en silice}\index{saturation en silice} s'il contient
plus de silice que ses autres oxydes ne peuvent en fixer dans des silicates, de sorte que du
quartz (ou un autre minéral de silice) apparaît~; \emph{sous-saturé en silice} s'il en contient
trop peu pour transformer tous ses alcalins en feldspath, de sorte qu'un feldspathoïde (néphéline,
leucite) apparaît à la place, jamais avec le quartz~; et \emph{saturé en silice}
entre les deux.
\end{definition}

Pour une lave dont les minéraux ne peuvent être comptés, la saturation se juge sur une
\emph{norme}~: un calcul standard qui répartit les oxydes de l'analyse entre des minéraux idéaux,
dans un ordre fixé, les alcalins d'abord dans le feldspath, et qui donne l'excès de silice
sous forme de quartz ou le déficit sous forme de feldspathoïde. La norme la plus usuelle, la norme CIPW, sépare
les deux familles de basaltes~: un basalte à néphéline dans sa norme est un basalte alcalin~;
un basalte sans néphéline est sub-alcalin, et parmi eux les basaltes tholéiitiques portent de
l'hypersthène (orthopyroxène) dans la norme et, souvent, un pyroxène pauvre en calcium dans le mode.
\Cref{exo:b2:magmatic-series:11} propose une version simplifiée de ce calcul.
% ledger: bgs:tas-use, bgs:series-terms

\section{Séries magmatiques}\label{sec:b2:magmatic-series:series}

Les laves d'un même volcan, ou d'une même province, se ressemblent rarement, mais elles sont apparentées~: elles
dérivent de \omterm{def:b2:magma-emplacement:magma}{magmas} parents voisins par la différenciation étudiée au
\cref{ch:b2:magmatic-differentiation}, et elles tracent ensemble une même ligne de descente.

\begin{definition}[Série magmatique]\label{def:b2:magmatic-series:series}
Une \emph{série magmatique}\index{série magmatique} est un ensemble de roches magmatiques d'un même cycle
ou d'une même province, apparentées par différenciation à des \omterm{def:b2:magma-emplacement:magma}{magmas} parents semblables et partageant des
caractères chimiques et minéralogiques propres. On en distingue couramment trois. La
\emph{série tholéiitique}\index{série tholéiitique} part de basaltes saturés ou légèrement
\omterm{def:b2:magmatic-series:saturation}{sursaturés en silice}, avec de l'hypersthène dans la norme, et montre un fort enrichissement en fer
au début de sa différenciation. La \emph{série calco-alcaline}\index{série calco-alcaline} va du
basalte à l'andésite, puis à la dacite et à la rhyolite, avec peu ou pas d'enrichissement en fer, et
avec des minéraux hydratés (hornblende, biotite) dans ses termes évolués. La
\emph{série alcaline}\index{série alcaline} part de basaltes à néphéline normative et évolue, à travers
les champs S et U du diagramme TAS,
vers le trachyte et la phonolite.
\end{definition}
% ledger: bgs:series-terms, bgs:tas-use

Ces noms désignent des ensembles de roches, et non des roches isolées~: un basalte isolé peut
appartenir à l'un quelconque d'entre eux, et les mêmes noms ont reçu des sens différents, si bien
qu'une pétrologue doit préciser ce qu'elle entend par eux. La distinction entre les \omterm{def:b2:magmatic-series:series}{séries
tholéiitique} et calco-alcaline repose sur le fer, et le diagramme qui la montre le mieux porte
sur trois groupes d'oxydes.
% ledger: bgs:series-terms

\begin{definition}[Diagramme AFM]\label{def:b2:magmatic-series:afm}
Un \emph{diagramme AFM}\index{diagramme AFM} est un diagramme triangulaire dont les sommets sont
$A = \ce{Na2O + K2O}$, $F = \ce{FeO_T}$ (tout le fer en \ce{FeO}) et $M = \ce{MgO}$, en
pourcentage massique recalculé de sorte que $A + F + M = 100$. Une droite tracée dessus sépare le
champ tholéiitique, du côté de F, du champ calco-alcalin.
\end{definition}
% ledger: vp21:series, geoplotr:afm

Un basalte en cours de différenciation perd du magnésium plus vite que tout autre élément,
puisque l'olivine et le pyroxène sont plus riches en magnésium que le liquide, et il garde ses
alcalins, qui sont incompatibles~: sa ligne de descente s'éloigne de M, vers A. Ce que devient le
fer décide de la série. Dans un \omterm{def:b2:magma-emplacement:magma}{magma} tholéiitique, le fer ne part qu'avec les silicates ferreux
riches en magnésium, plus lentement que lui, de sorte que le rapport $F/M$ augmente et que la
trajectoire monte d'abord vers F avant de tourner vers A~: c'est la \emph{tendance de Fenner}.
Dans un \omterm{def:b2:magma-emplacement:magma}{magma} calco-alcalin, plus oxydé, le fer part aussi dans la magnétite, qui contient du fer
ferrique, et $F/M$ reste presque constant pendant que le liquide se dirige en ligne droite vers A~:
c'est la \emph{tendance de Bowen}. La figure ci-dessous montre les deux avec de vraies roches,
les tholéiites d'Islande et les laves calco-alcalines des Cascades. La ligne de partage, tracée à
l'œil en 1971, les sépare bien là où elles sont basaltiques~; à fortes teneurs en alcalins, les deux
trajectoires convergent et la ligne en dit peu.
% ledger: vp21:series, geoplotr:afm, data:b4-cath

\begin{omfigure}
\centering
\begin{tikzpicture}
\begin{axis}[name=afm, width=0.5\linewidth, height=0.5\linewidth, axis equal image,
  hide axis, xmin=-0.04, xmax=1.04, ymin=-0.06, ymax=0.92, clip=false]
\addplot[only marks, mark=*, mark size=0.5pt, blue!70!black, opacity=0.6] table[x=afm_x, y=afm_y] {figdata/out/bachelor-2-series-cath-th.dat};
\addplot[only marks, mark=*, mark size=0.5pt, red!75!black, opacity=0.6] table[x=afm_x, y=afm_y] {figdata/out/bachelor-2-series-cath-ca.dat};
\addplot[very thick, black] table[x=x, y=y] {figdata/out/bachelor-2-series-afm-line.dat};
\draw[thick] (axis cs:0,0) -- (axis cs:1,0) -- (axis cs:0.5,0.8660254) -- cycle;
\node[anchor=south] at (axis cs:0.5,0.866) {F};
\node[anchor=north east] at (axis cs:0,0) {A};
\node[anchor=north west] at (axis cs:1,0) {M};
\node[font=\scriptsize, fill=white, inner sep=1pt] at (axis cs:0.52,0.62) {tholéiitique};
\node[font=\scriptsize, fill=white, inner sep=1pt] at (axis cs:0.33,0.12) {calco-alcalin};
\end{axis}
\begin{axis}[omrecord, at={(afm.east)}, anchor=west, xshift=1.3cm,
  width=0.46\linewidth, height=0.46\linewidth,
  xmin=-2.5, xmax=2.2, ymin=-4.5, ymax=1.5,
  xlabel={$\ln(A/F)$}, ylabel={$\ln(M/F)$},
  legend style={font=\scriptsize, at={(0.98,0.98)}, anchor=north east},
  tick label style={font=\scriptsize}, label style={font=\small}]
\addplot[only marks, mark=*, mark size=0.5pt, blue!70!black, opacity=0.6] table[x=lnAF, y=lnMF] {figdata/out/bachelor-2-series-lr-th.dat};
\addlegendentry{Islande}
\addplot[only marks, mark=*, mark size=0.5pt, red!75!black, opacity=0.6] table[x=lnAF, y=lnMF] {figdata/out/bachelor-2-series-lr-ca.dat};
\addlegendentry{Cascades}
\addplot[thick, blue!50!black, forget plot] table[x=x, y=y_th] {figdata/out/bachelor-2-series-lr-fits.dat};
\addplot[thick, red!50!black, forget plot] table[x=x, y=y_ca] {figdata/out/bachelor-2-series-lr-fits.dat};
\end{axis}
\end{tikzpicture}
\par\smallskip{\small À gauche~: le \omterm{def:b2:magmatic-series:afm}{diagramme AFM} des roches islandaises (bleu) et des Cascades (rouge)
de la figure TAS, avec la ligne de partage d'Irvine et Baragar (1971). À droite~: les mêmes roches
en coordonnées de rapports logarithmiques, où chaque série devient une bande presque droite~;
les droites sont des régressions orthogonales, de pentes $-1.10$ (Islande) et $-0.72$ (Cascades),
se rejoignant en $(-3.84, 2.12)$ comme l'ont trouvé Vermeesch et Pease (2021).}
\end{omfigure}
% ledger: geoplotr:afm, data:b4-cath, vp21:model

Parmi les roches islandaises situées dans le domaine de la ligne de partage, 351 tombent du côté
tholéiitique et 23 de l'autre~; parmi les roches des Cascades, 417 tombent du côté calco-alcalin et 57
de l'autre. Une courbe est une frontière peu commode, et un modèle simple explique la forme des
tendances. Supposons qu'un \omterm{def:b2:magma-emplacement:magma}{magma} perde ses alcalins, son fer et son magnésium à des vitesses
proportionnelles aux quantités qu'il contient, $\mathrm dX/\mathrm dt = -\lambda_X X$.

\begin{proposition}[Tendances en rapports logarithmiques]\label{prop:b2:magmatic-series:logratio}
Si $A$, $F$ et $M$ décroissent exponentiellement avec les constantes $\lambda_A$, $\lambda_F$ et
$\lambda_M$, le \omterm{def:b2:magma-emplacement:magma}{magma} se déplace le long d'une droite dans le plan de $\ln(A/F)$ et de
$\ln(M/F)$, de pente $(\lambda_F - \lambda_M)/(\lambda_F - \lambda_A)$, alors qu'il suit une courbe
sur le triangle AFM.
\end{proposition}

\begin{proof}
Avec $X = X_0 e^{-\lambda_X t}$, $\ln(A/F) = \ln(A_0/F_0) + (\lambda_F - \lambda_A)t$
et $\ln(M/F) = \ln(M_0/F_0) + (\lambda_F - \lambda_M)t$~: toutes deux sont linéaires en $t$, donc
l'une est linéaire en l'autre, de la pente annoncée. Les rapports $A/F$ et $M/F$ fixent le
point sur le triangle, dont les coordonnées $A/(A + F + M)$, \dots\ ne sont pas linéaires en $t$.
\end{proof}

La pente plus raide des roches islandaises signifie un $\lambda_F$ plus petit~: le fer quitte
le \omterm{def:b2:magma-emplacement:magma}{magma} tholéiitique plus lentement (\cref{exo:b2:magmatic-series:10}). Vermeesch et Pease
ont montré que les deux bandes de la figure ci-dessus se rejoignent en un seul point, qu'ils
lisent comme une source commune hypothétique.
% ledger: vp21:model

Les \omterm{def:b2:magmatic-series:series}{séries tholéiitique} et calco-alcaline diffèrent aussi par leur teneur en potasse. Le diagramme
\ce{K2O}--\ce{SiO2} répartit les basaltes, andésites, dacites et rhyolites à plus de \qty{48}{\percent} de silice
en types bas-K, moyen-K et haut-K, par deux droites (la figure ci-dessous). La plupart des deux
suites étudiées ici sont de type moyen-K~: le diagramme indique combien de potasse une série
porte, et non de quelle série il s'agit.
% ledger: bgs:k2o-sio2, data:b4-cath

\begin{omfigure}
\centering
\begin{tikzpicture}
\begin{axis}[omrecord, width=0.78\linewidth, height=6.2cm,
  xmin=45, xmax=79, ymin=0, ymax=5,
  xlabel={\ce{SiO2} (\unit{\percent})}, ylabel={\ce{K2O} (\unit{\percent})},
  legend style={font=\scriptsize, at={(0.02,0.98)}, anchor=north west},
  tick label style={font=\scriptsize}, label style={font=\small}]
\addplot[only marks, mark=*, mark size=0.55pt, blue!70!black, opacity=0.6] table[x=SiO2, y=K2O] {figdata/out/bachelor-2-series-cath-th.dat};
\addlegendentry{Islande}
\addplot[only marks, mark=*, mark size=0.55pt, red!75!black, opacity=0.6] table[x=SiO2, y=K2O] {figdata/out/bachelor-2-series-cath-ca.dat};
\addlegendentry{Cascades}
\addplot[thick, black, forget plot] table[x=x, y=y_low] {figdata/out/bachelor-2-series-k2o-lines.dat};
\addplot[thick, black, forget plot] table[x=x, y=y_high] {figdata/out/bachelor-2-series-k2o-lines.dat};
\draw[thick] (axis cs:48,0) -- (axis cs:48,0.3);
\node[font=\scriptsize, fill=white, inner sep=1pt] at (axis cs:60,3.7) {haut-K};
\node[font=\scriptsize, fill=white, inner sep=1pt] at (axis cs:76.5,2.2) {moyen-K};
\node[font=\scriptsize, fill=white, inner sep=1pt] at (axis cs:60,0.35) {bas-K};
\end{axis}
\end{tikzpicture}
\par\smallskip{\small Teneur en potasse en fonction de la silice pour les roches islandaises et des Cascades, avec
les droites passant par (48, 0.3)--(68, 1.2) et (48, 1.2)--(68, 3.1) qui séparent les types bas-K,
moyen-K et haut-K (d'après Peccerillo et Taylor, 1976, tel qu'adopté par le schéma international),
prolongées ici au-delà de 68.}
\end{omfigure}
% ledger: bgs:k2o-sio2, data:b4-cath

Pourquoi l'oxydation diffère-t-elle~? Les \omterm{def:b2:magma-emplacement:magma}{magmas} calco-alcalins se forment au-dessus des zones de subduction, où
la plaque plongeante, en se déshydratant, apporte au coin mantellique des espèces oxydées~; ces \omterm{def:b2:magma-emplacement:magma}{magmas}
cristallisent de la magnétite, et leurs termes évolués de la hornblende ou de la biotite. Les \omterm{def:b2:magma-emplacement:magma}{magmas} tholéiitiques
naissent sur les dorsales, où le manteau fond par décompression
(\cref{ch:b2:mantle-melting}), et sont plus réduits.
% ledger: vp21:series

\section{Séries et contextes}\label{sec:b2:magmatic-series:settings}

Chaque contexte de la tectonique des plaques produit ses propres \omterm{def:b2:magma-emplacement:magma}{magmas}, car chacun fait fondre
le manteau à sa manière~: par décompression sous les dorsales et les points chauds, par apport d'eau
au-dessus des zones de subduction (\cref{ch:b2:mantle-melting}).

Les \emph{dorsales océaniques} émettent des basaltes tholéiitiques, les MORB du
\cref{ch:b2:elements-in-the-earth}, pauvres en \omterm{def:b2:elements-in-the-earth:partition}{éléments incompatibles} et à signature isotopique
appauvrie de leur source. L'Islande, dorsale au-dessus d'un point chaud, ajoute à cette liste ses
propres tholéiites, avec une suite complète jusqu'à la rhyolite.
% ledger: vp21:series

Les \emph{îles océaniques} croissent par étapes. Un volcan hawaïen émet des laves alcalines pendant sa
jeunesse, puis édifie son bouclier de basalte tholéiitique, estimé entre 80 et
\qty{95}{\percent} de son volume final, au cours de ses 1 à 2 premiers Ma~; en s'éloignant du
point chaud, ses laves redeviennent alcalines, et finissent par une longue traîne d'éruptions
alcalines rares. Une fusion moindre donne des \omterm{def:b2:magma-emplacement:magma}{magmas} plus alcalins~: les alcalins, incompatibles,
sont moins dilués dans un faible \omterm{def:b2:mantle-melting:fraction}{taux de fusion}. Du point de vue isotopique, les \omterm{def:b2:elements-in-the-earth:morb-oib}{basaltes d'île
océanique} échantillonnent les réservoirs enrichis du \cref{ch:b2:isotope-geochemistry}.
% ledger: usgs14:hawaii-stages

Les \emph{arcs volcaniques} émettent des \omterm{def:b2:magmatic-series:series}{séries calco-alcalines}. Leurs \omterm{def:b2:elements-in-the-earth:trace}{éléments en traces} portent
une marque distinctive (figure ci-dessous)~: normalisés au \omterm{def:b2:elements-in-the-earth:bse}{manteau primitif}, les basaltes d'arc
sont enrichis en Rb, Ba, Th, U, Pb et Sr, éléments que la plaque plongeante ajoute au coin
mantellique, et appauvris en niobium et en tantale par rapport à leurs voisins. Le basalte d'arc
médian de la base de données a $\mathrm{Nb}/\mathrm{La} = 0.49$, contre 0.97 pour les MORB
et 1.23 pour les \omterm{def:b2:elements-in-the-earth:morb-oib}{basaltes d'île océanique}.
% ledger: data:b4-basalt-settings

\begin{omfigure}
\centering
\begin{tikzpicture}
\begin{axis}[omrecord, width=0.86\linewidth, height=6.4cm, ymode=log,
  xmin=0.5, xmax=22.5, ymin=1, ymax=300,
  xtick={1,...,22}, xticklabels={Rb,Ba,Th,U,Nb,Ta,La,Ce,Pb,Sr,Nd,Zr,Hf,Sm,Eu,Ti,Gd,Dy,Y,Er,Yb,Lu},
  ylabel={échantillon / manteau primitif},
  legend style={font=\scriptsize, at={(0.98,0.98)}, anchor=north east},
  tick label style={font=\scriptsize}, label style={font=\small}]
\addplot[thick, mark=square*, mark size=1.4pt, orange!85!black] table[x=idx, y=oib] {figdata/out/bachelor-2-series-spider.dat};
\addlegendentry{basalte d'île océanique}
\addplot[thick, mark=*, mark size=1.4pt, blue!70!black] table[x=idx, y=morb] {figdata/out/bachelor-2-series-spider.dat};
\addlegendentry{MORB}
\addplot[thick, mark=triangle*, mark size=1.7pt, red!75!black] table[x=idx, y=iab] {figdata/out/bachelor-2-series-spider.dat};
\addlegendentry{basalte d'arc insulaire}
\end{axis}
\end{tikzpicture}
\par\smallskip{\small Compositions médianes normalisées au \omterm{def:b2:elements-in-the-earth:bse}{manteau primitif} de 241 MORB, 259
\omterm{def:b2:elements-in-the-earth:morb-oib}{basaltes d'île océanique} et 256 basaltes d'arc insulaire (la compilation de Vermeesch, 2006~;
composition du \omterm{def:b2:elements-in-the-earth:bse}{manteau primitif} de McDonough et Sun, 1995). Les basaltes d'arc présentent le creux en Nb
et Ta et des pics en Ba, Pb et Sr~; les \omterm{def:b2:elements-in-the-earth:morb-oib}{basaltes d'île océanique} sont enrichis dans tous les
éléments les plus incompatibles.}
\end{omfigure}
% ledger: data:b4-basalt-settings, ms95:pm-trace

Dans les continents, le \omterm{def:b2:magma-emplacement:magma}{magma} basaltique rencontre une croûte épaisse qu'il peut chauffer et fondre,
et le résultat peut être un volcanisme à deux modes.

\begin{definition}[Volcanisme bimodal]\label{def:b2:magmatic-series:bimodal}
Le \emph{volcanisme bimodal}\index{volcanisme bimodal} est l'éruption, dans une même province,
de deux groupes de \omterm{def:b2:magma-emplacement:magma}{magmas}, basaltiques et rhyolitiques, avec peu de compositions intermédiaires
entre eux.
\end{definition}

Le volcanisme le long de la trace du point chaud de Yellowstone est bimodal~: sur la plaine
orientale du Snake River, des rhyolites plus anciennes, à environ \qty{75}{\percent} de silice, sont
recouvertes de basaltes plus jeunes, à environ \qty{50}{\percent}. L'écart entre les deux est l'indice~:
la différenciation d'un seul basalte aurait laissé des laves intermédiaires en place, et un écart
signale une seconde source pour les rhyolites, comme la croûte que le basalte réchauffe.
% ledger: yvo23:bimodal

Les \emph{grandes provinces ignées}, rencontrées dans le volume de première année, émettent des basaltes en volumes
qu'aucun autre contexte n'égale~; leur série se lit avec les mêmes diagrammes.

\section{Diagrammes de discrimination}\label{sec:b2:magmatic-series:discrimination}

Les régularités de la section précédente suggèrent de lire un contexte directement dans une analyse,
et les pétrologues ont tracé de nombreux diagrammes pour cela.

\begin{definition}[Diagramme de discrimination]\label{def:b2:magmatic-series:discrimination}
Un \emph{diagramme de discrimination}\index{diagramme de discrimination} est un diagramme de
variables chimiques, généralement des \omterm{def:b2:elements-in-the-earth:trace}{éléments en traces} immobiles, divisé en champs tracés autour des
compositions de roches de contextes connus, et servant à attribuer un contexte à une roche d'origine
inconnue.
\end{definition}

Le classique est le triangle Ti--Zr--Y de Pearce et Cann (1973), tracé pour les basaltes. Il
utilise des éléments peu déplacés par l'altération, si bien qu'il s'applique aux laves anciennes et
métamorphisées dont le contexte est perdu. La figure ci-dessous le teste sur les 538 basaltes de la
compilation qui possèdent les trois éléments. La plupart des \omterm{def:b2:elements-in-the-earth:morb-oib}{basaltes d'île océanique} tombent dans le champ
intraplaque D, 152 sur 173, et la plupart des MORB dans le champ B, 142 sur 164~; mais le champ B est
partagé avec les basaltes d'arc, et 101 des 201 basaltes d'arc y tombent aussi.
% ledger: geoplotr:tizry, data:b4-basalt-settings

\begin{omfigure}
\centering
\begin{tikzpicture}
\begin{axis}[width=0.62\linewidth, height=0.55\linewidth, axis equal image,
  hide axis, xmin=-0.04, xmax=1.04, ymin=-0.06, ymax=0.92, clip=false,
  legend style={font=\scriptsize, at={(1.0,0.95)}, anchor=north west}]
\addplot[only marks, mark=*, mark size=0.7pt, blue!70!black, opacity=0.7] table[x=x, y=y] {figdata/out/bachelor-2-series-tizry-morb.dat};
\addlegendentry{MORB}
\addplot[only marks, mark=square*, mark size=0.7pt, orange!85!black, opacity=0.7] table[x=x, y=y] {figdata/out/bachelor-2-series-tizry-oib.dat};
\addlegendentry{basalte d'île océanique}
\addplot[only marks, mark=triangle*, mark size=0.9pt, red!75!black, opacity=0.7] table[x=x, y=y] {figdata/out/bachelor-2-series-tizry-iab.dat};
\addlegendentry{basalte d'arc insulaire}
\foreach \f in {A,B,C,D}
  \addplot[thick, black, forget plot] table[x=x, y=y] {figdata/out/bachelor-2-series-tizry-\f.dat};
\draw[thick] (axis cs:0,0) -- (axis cs:1,0) -- (axis cs:0.5,0.8660254) -- cycle;
\node[anchor=south] at (axis cs:0.5,0.866) {Ti/100};
\node[anchor=north east] at (axis cs:0,0) {Zr};
\node[anchor=north west] at (axis cs:1,0) {Y $\times$ 3};
\draw[thin] (axis cs:0.33,0.36) -- (axis cs:0.22,0.46) node[anchor=south east, font=\small\bfseries, inner sep=1pt] {D};
\draw[thin] (axis cs:0.36,0.12) -- (axis cs:0.24,0.06) node[anchor=north east, font=\small\bfseries, inner sep=1pt] {C};
\draw[thin] (axis cs:0.56,0.37) -- (axis cs:0.76,0.50) node[anchor=west, font=\small\bfseries, inner sep=1pt] {A};
\draw[thin] (axis cs:0.53,0.21) -- (axis cs:0.72,0.13) node[anchor=north west, font=\small\bfseries, inner sep=1pt] {B};
\end{axis}
\end{tikzpicture}
\par\smallskip{\small Le diagramme Ti--Zr--Y de Pearce et Cann (1973) avec les basaltes
de contexte connu. A tholéiites d'arc insulaire~; B MORB, tholéiites d'arc insulaire et
basaltes calco-alcalins~; C basaltes calco-alcalins~; D basaltes intraplaques. Ti issu de
\ce{TiO2}~; les trois éléments en ppm.}
\end{omfigure}
% ledger: geoplotr:tizry, data:b4-basalt-settings

\begin{method}[Utiliser un diagramme de discrimination]\label{met:b2:magmatic-series:discrimination}
(1) Vérifier que la roche est du type pour lequel le diagramme a été construit, en général un
basalte, et non un \omterm{def:b2:magmatic-differentiation:cumulate}{cumulat} ni une lave évoluée, dont les \omterm{def:b2:elements-in-the-earth:trace}{éléments en traces} reflètent autant la
cristallisation que la source. (2) Préférer, pour les roches altérées ou métamorphisées, des
éléments immobiles dans l'altération (Ti, Zr, Y, Nb, les \omterm{def:b2:elements-in-the-earth:ree}{terres rares} lourdes). (3) Tracer une suite
d'échantillons, et non un seul~: un point isolé près d'une limite ne décide de rien. (4) Lire le
résultat comme une vraisemblance~: les champs ont été tracés autour de la plupart, et non de la
totalité, de leurs échantillons de référence, et certains champs sont partagés. (5) Confirmer par des
preuves indépendantes~: d'autres diagrammes, les isotopes, le contexte géologique.
\end{method}

L'étape (4) peut être rendue quantitative (\cref{exo:b2:magmatic-series:8}). Si un basalte tombe dans
le champ B, et si les trois contextes étaient également fréquents, la probabilité qu'il s'agisse
d'un MORB n'est que d'environ 0.62. Les classifieurs statistiques entraînés sur de telles compilations
héritent de la même limite~: un basalte porte la mémoire de sa source et de sa fusion, et deux
contextes peuvent produire des \omterm{def:b2:magma-emplacement:magma}{magmas} semblables.
% ledger: data:b4-basalt-settings

\begin{inthelab}[Contrôler une analyse de roche totale]
Avant de placer une roche sur l'un de ces diagrammes, les pétrologues contrôlent l'analyse. Les
oxydes majeurs, mesurés par exemple par fluorescence X, doivent totaliser une valeur voisine de 100
avec les volatils~: l'eau et le dioxyde de carbone, mesurés directement ou ensemble sous forme de
perte au feu, c'est-à-dire la masse qu'une poudre perd sous un chauffage intense. Une forte
perte au feu signale une roche altérée, dont les alcalins ont pu se déplacer~; le schéma du TAS rejette
les analyses contenant plus de \qty{2}{\percent} d'eau ou
\qty{0.5}{\percent} de dioxyde de carbone. Les analyses sont ensuite recalculées à 100 sans
volatils, le fer est exprimé en \ce{FeO} total ou réparti entre fer ferreux et ferrique, et
les matériaux de référence analysés dans la même série sont comparés à leurs valeurs certifiées
(\cref{ch:b2:field-lab-methods-2}).
\end{inthelab}
% ledger: bgs:tas-use

\section{Exercices}

\begin{exercise}[$\star$]\label{exo:b2:magmatic-series:1}
Nommer les roches volcaniques suivantes sur le diagramme TAS (\ce{SiO2}, \ce{Na2O}, \ce{K2O}
en \unit{\percent}, déjà exemptes de volatils)~: (a) 49.60, 2.22, 0.51 (le basalte de référence
hawaïen BHVO-2)~; (b) 58.5, 3.9, 1.6~; (c) 72.0, 4.0, 4.2~; (d) 47.0, 4.5, 2.5.
Pour (d), nommer les deux roches possibles et dire ce qui décide entre elles.
\end{exercise}
% ledger: usgs:bhvo2-major, bgs:tas

\begin{exercise}[$\star$]\label{exo:b2:magmatic-series:2}
Une lave donne \ce{SiO2} 51.60, \ce{TiO2} 1.20, \ce{Al2O3} 16.10, \ce{FeO_T} 9.80,
\ce{MgO} 6.60, \ce{CaO} 8.20, \ce{Na2O} 3.50, \ce{K2O} 1.40, \ce{H2O} 1.40 et
\ce{CO2} \qty{0.20}{\percent}. (a) L'analyse est-elle acceptable pour le diagramme TAS~?
(b) La recalculer exempte de volatils. (c) Nommer la roche avant et après le
recalcul. (d) Est-elle de type bas-K, moyen-K ou haut-K~?
\end{exercise}
% ledger: bgs:tas-use, bgs:k2o-sio2

\begin{exercise}[$\star$]\label{exo:b2:magmatic-series:3}
Nommer sur le \omterm{def:b2:magmatic-series:qapf}{diagramme QAPF}~: (a) une roche grenue à \qty{25}{\percent} de quartz,
\qty{30}{\percent} de feldspath alcalin, \qty{35}{\percent} de plagioclase, \qty{8}{\percent}
de biotite et \qty{2}{\percent} de hornblende~; (b) une roche à \qty{5}{\percent} de quartz,
\qty{5}{\percent} de feldspath alcalin, \qty{55}{\percent} de plagioclase et
\qty{35}{\percent} de clinopyroxène.
\end{exercise}
% ledger: bgs:qapf

\begin{exercise}[$\star$]\label{exo:b2:magmatic-series:4}
Classer en type bas-K, moyen-K ou haut-K, à l'aide des droites de la figure \ce{K2O}--\ce{SiO2}~:
(a) 50.0, 0.40~; (b) 56.0, 1.50~; (c) 62.0, 2.90~; (d) 66.0, 1.00 (\ce{SiO2},
\ce{K2O}, \unit{\percent}).
\end{exercise}
% ledger: bgs:k2o-sio2

\begin{exercise}[$\star\star$]\label{exo:b2:magmatic-series:5}
(a) Calculer $A$, $F$ et $M$ pour un basalte à \ce{FeO_T} 14.5, \ce{MgO} 4.5 et
\ce{Na2O + K2O} \qty{3.2}{\percent}, et pour une andésite à 6.1, 3.8 et 5.05.
(b) La ligne de partage passe par $(F, A, M) = (54.5, 23.0, 22.5)$, $(52.5, 30.0,
17.5)$ et $(49.6, 35.0, 15.4)$, le premier point étant le plus riche en fer~; le
côté tholéiitique est celui de F. De quel côté tombe chaque roche~? (c) Quel
minéral, cristallisant tôt, maintient l'andésite à sa place~?
\end{exercise}
% ledger: geoplotr:afm, vp21:series

\begin{exercise}[$\star\star$]\label{exo:b2:magmatic-series:6}
Le MORB médian de la compilation contient \qty{4}{ppm} de Nb et \qty{4.12}{ppm} de La, le
\omterm{def:b2:elements-in-the-earth:morb-oib}{basalte d'île océanique} médian 38 et 30.8, le basalte d'arc médian 3 et 6.1. Avec les valeurs du
\omterm{def:b2:elements-in-the-earth:bse}{manteau primitif} Nb 0.658 et La \qty{0.648}{ppm}, calculer les valeurs normalisées
et $\mathrm{Nb_N}/\mathrm{La_N}$ pour chacun. Lequel présente une anomalie en niobium, et pourquoi~?
\end{exercise}
% ledger: data:b4-basalt-settings, ms95:pm-trace

\begin{exercise}[$\star\star$]\label{exo:b2:magmatic-series:7}
La masse de Ti vaut $0.5993$ fois celle de \ce{TiO2}. Calculer les coordonnées dans le diagramme Ti--Zr--Y
(Ti/100, Zr, $3\mathrm Y$, ramenées à 100) de trois basaltes~: (a) \ce{TiO2}
\qty{1.6}{\percent}, Zr 110, Y \qty{32}{ppm}~; (b) 2.8, 220, 29~; (c) 0.8, 50, 20. Les placer
sur la figure et dire ce que chaque placement permet de conclure.
\end{exercise}
% ledger: geoplotr:tizry

\begin{exercise}[$\star\star$]\label{exo:b2:magmatic-series:8}
Dans le test Ti--Zr--Y de ce chapitre, le champ B contient 142 MORB sur 164, 4 \omterm{def:b2:elements-in-the-earth:morb-oib}{basaltes d'île
océanique} sur 173 et 101 basaltes d'arc sur 201. (a) Quelle fraction de chaque groupe tombe
dans B~? (b) Si les trois contextes étaient également fréquents, quelle est la probabilité qu'un basalte
tombant dans B soit un MORB~? (c) Que permet d'affirmer le seul champ B~?
\end{exercise}
% ledger: data:b4-basalt-settings

\begin{exercise}[$\star\star$]\label{exo:b2:magmatic-series:9}
Une rhyolite contient \ce{Na2O} 5.5, \ce{K2O} 4.5 et \ce{Al2O3} \qty{10.5}{\percent} ($M$ =
61.98, 94.20 et \qty{101.96}{g/mol}). Calculer son rapport molaire $(\ce{Na2O +
K2O})/\ce{Al2O3}$. (a) Est-elle hyperalcaline~? (b) Le feldspath alcalin fixe un oxyde
alcalin par \ce{Al2O3}~: quelle fraction des alcalins ne peut y entrer~?
\end{exercise}
% ledger: bgs:tas-use

\begin{exercise}[$\star\star\star$]\label{exo:b2:magmatic-series:10}
Dans le modèle de la proposition sur les rapports logarithmiques, prendre $\lambda_A = 0.4$ et $\lambda_M =
0.8$. (a) Calculer la pente de la tendance pour $\lambda_F = 0.5$ et pour $\lambda_F =
0.7$. (b) Les roches islandaises ont une pente de $-1.10$ et celles des Cascades
$-0.72$~: quel $\lambda_F$ chacune implique-t-elle~? (c) Quelle série perd son fer le plus vite,
et quel minéral l'explique~? (d) Que devient $F/M$ le long de chaque tendance~?
\end{exercise}
% ledger: vp21:model

\begin{exercise}[$\star\star\star$]\label{exo:b2:magmatic-series:11}
Une norme simplifiée. Travailler en moles d'oxyde. Mettre tout le \ce{K2O} dans l'orthose
(\ce{K2O.Al2O3.6SiO2}), tout le \ce{Na2O} dans l'albite (\ce{Na2O.Al2O3.6SiO2}), le reste de
\ce{Al2O3} avec autant de \ce{CaO} dans l'anorthite (\ce{CaO.Al2O3.2SiO2}), le reste de
\ce{CaO} avec autant de $(\ce{Mg,Fe})\ce{O}$ dans le diopside
(\ce{CaO.(Mg,Fe)O.2SiO2}), et le reste de $(\ce{Mg,Fe})\ce{O}$ dans l'hypersthène
(\ce{(Mg,Fe)O.SiO2}). Si la silice manque, transformer l'hypersthène en olivine
(\ce{2(Mg,Fe)O.SiO2})~; s'il en manque encore, transformer l'albite en néphéline
(\ce{Na2O.Al2O3.2SiO2}). Appliquer le calcul à (a) BHVO-2~: \ce{SiO2} 49.60, \ce{Al2O3} 13.44,
\ce{FeO_T} 11.15, \ce{MgO} 7.257, \ce{CaO} 11.40, \ce{Na2O} 2.219, \ce{K2O}
\qty{0.513}{\percent}~; (b) un basalte à 45.5, 15.0, 11.5, 8.5, 10.0, 3.8 et 1.5 (masses molaires
\ce{SiO2} 60.08, \ce{Al2O3} 101.96, \ce{FeO} 71.84, \ce{MgO} 40.30, \ce{CaO}
56.08, \ce{Na2O} 61.98, \ce{K2O} \qty{94.20}{g/mol}). Chacun est-il saturé en silice~? À quelle
série appartient chaque basalte~?
\end{exercise}
% ledger: usgs:bhvo2-major, bgs:series-terms

\begin{exercise}[$\star\star\star$]\label{exo:b2:magmatic-series:12}
Dans la compilation, $\mathrm{Nb}/\mathrm{La} < 0.7$ vaut pour 116 basaltes d'arc sur 141,
21 MORB sur 101 et 3 \omterm{def:b2:elements-in-the-earth:morb-oib}{basaltes d'île océanique} sur 137 analysés pour les deux. (a) Calculer les
trois fractions. (b) Avec des proportions égales des trois contextes, quelle est la
probabilité qu'un basalte de $\mathrm{Nb}/\mathrm{La} < 0.7$ soit un basalte d'arc~? (c)
Refaire le calcul pour une région où \qty{60}{\percent} des basaltes sont des MORB, \qty{30}{\percent}
des \omterm{def:b2:elements-in-the-earth:morb-oib}{basaltes d'île océanique} et \qty{10}{\percent} des basaltes d'arc. (d) Quelle leçon le (c)
apporte-t-il à l'usage des \omterm{def:b2:magmatic-series:discrimination}{diagrammes de discrimination}~?
\end{exercise}
% ledger: data:b4-basalt-settings

\section{Problème~: un volcan, trois laves}

\begin{problem}[{Problème du week-end --- de trois analyses au contexte d'un volcan
et aux cristaux qu'il a laissés derrière lui}]
\label{pb:b2:magmatic-series:1}
Trois laves d'un stratovolcan (fictif), de ses plus anciennes coulées de flanc à son dôme le plus
récent, donnent (\unit{\percent})~:

\smallskip\noindent{\small
\begin{tabular}{@{}lrrrrrrrrrr@{}}
 & \ce{SiO2} & \ce{TiO2} & \ce{Al2O3} & \ce{FeO_T} & \ce{MgO} & \ce{CaO} & \ce{Na2O} & \ce{K2O} & \ce{H2O} & \ce{CO2}\\
lave 1 & 50.4 & 0.95 & 17.6 & 8.6 & 7.9 & 10.2 & 2.9 & 0.62 & 0.70 & 0.10\\
lave 2 & 58.6 & 0.80 & 17.4 & 6.1 & 3.8 & 6.9 & 3.9 & 1.15 & 1.10 & 0.05\\
lave 3 & 65.3 & 0.55 & 16.2 & 4.0 & 1.6 & 4.3 & 4.4 & 1.85 & 1.20 & 0.05\\
\end{tabular}}

\smallskip\noindent et en ppm~: Rb 10, 20, 32~; Th 1.1, 2.1, 3.4~; Nb 2.0, 3.6, 5.6~;
La 7, 13, 20~; Sr 520, 480, 380~; Zr 70, 125, 190~; Y 18, 22, 24. \omterm{def:b2:elements-in-the-earth:bse}{Manteau primitif}~: Nb
0.658, La 0.648, Th \qty{0.0795}{ppm}.

\noindent\textbf{Partie I --- Nommer les laves.}
\begin{enumerate}
\item Additionner chaque analyse. Les trois sont-elles acceptables pour le diagramme TAS~?
\item Recalculer la lave 1 exempte de volatils~: donner le facteur, \ce{SiO2} et
\ce{Na2O + K2O}.
\item Faire de même pour les laves 2 et 3.
\item Nommer les trois laves sur le diagramme TAS.
\item Les noms changeraient-ils sans le recalcul~?
\item Une roche grenue de la composition de la lave 2 a cristallisé dans les racines
du volcan. Quel schéma la nomme, et dans quel champ l'attendez-vous~?
\end{enumerate}

\noindent\textbf{Partie II --- Quelle série~?}
\begin{enumerate}[resume]
\item Classer chaque lave en type bas-K, moyen-K ou haut-K.
\item Calculer $A$, $F$, $M$ pour chacune (exemptes de volatils).
\item La ligne de partage passe par $(F, A, M) = (42.8, 17.1, 40.1)$, $(46.9, 18.1,
35.0)$, $(52.5, 30.0, 17.5)$, $(49.6, 35.0, 15.4)$, $(40.8, 50.0, 9.2)$ et $(37.9,
55.0, 7.1)$, le côté tholéiitique étant celui de F. En interpolant $F$ linéairement en
$A$ entre points voisins, de quel côté tombe chaque lave~?
\item Calculer \ce{FeO_T}/\ce{MgO} pour chaque lave. Y a-t-il un enrichissement en fer~?
\item À quelle série appartient le volcan~? Quel minéral et quelles conditions
expliquent-elles cela~?
\end{enumerate}

\noindent\textbf{Partie III --- Quel contexte~?}
\begin{enumerate}[resume]
\item Normaliser Nb, La et Th de la lave 1 au \omterm{def:b2:elements-in-the-earth:bse}{manteau primitif}.
\item Calculer $\mathrm{Nb_N}/\mathrm{La_N}$ pour la lave 1 et la comparer aux médianes
de l'\cref{exo:b2:magmatic-series:6}.
\item Calculer $\mathrm{Th_N}/\mathrm{Nb_N}$. Qu'est-ce qui a transporté le thorium~?
\item Calculer les coordonnées Ti--Zr--Y de la lave 1 (\ce{TiO2} de l'analyse exempte de
volatils). Que dit son champ~? Pourquoi ne pas placer les laves 2 et 3~?
\item Conclure~: où ce volcan s'est-il édifié~?
\end{enumerate}

\noindent\textbf{Partie IV --- Quelle quantité a cristallisé~?}
\begin{enumerate}[resume]
\item La potasse n'entre dans aucun des premiers minéraux. À partir du \ce{K2O} exempt de
volatils des laves 1 et 2, calculer la fraction $F$ de liquide restant lorsque la lave 2 s'est formée, si elle
dérive de la lave 1 par \omterm{def:b2:magmatic-differentiation:crystallisation}{cristallisation fractionnée}.
\item Faire de même avec le rubidium et le thorium.
\item Le zirconium a un $D \approx 0.1$ global. Utiliser la loi de Rayleigh pour estimer $F$ à partir
de cet élément.
\item Le strontium diminue de la lave 2 à la lave 3. Calculer $F$ à partir de la potasse entre les laves 2 et
3, et s'en servir pour estimer le $D$ global du Sr sur cet intervalle. Quel minéral en est
responsable~?
\item De la lave 1 à la lave 3, quelle fraction du liquide a cristallisé~?
\item Si le \omterm{def:b2:magma-emplacement:magma}{magma} avait assimilé une croûte granitique riche en potasse, l'estimation de la potasse
serait-elle trop haute ou trop basse~?
\item Faire la moyenne des quatre estimations de la partie IV pour les laves 1 à 2 et donner la masse de
cristaux laissés par \qty{100}{kg} de basalte pour produire l'andésite, le nombre clé
de ce problème.
\end{enumerate}
\end{problem}
% ledger: bgs:tas-use, bgs:tas, bgs:k2o-sio2, geoplotr:afm, geoplotr:tizry, ms95:pm-trace
